\documentclass[10pt,a4paper]{amsart}
\usepackage[margin=19mm]{geometry}
\usepackage{amsmath,amssymb,mathtools}
\usepackage[scr=rsfs,cal=euler]{mathalfa}
\usepackage{xfrac}
\usepackage[unicode,hidelinks,bookmarks=false]{hyperref}
\newcommand{\ie}{i.e.\ }
\newcommand{\cf}{cf.\ }
\newcommand{\resp}{respectively\ }
\newcommand{\Subsection}{\subsection}
\providecommand{\nobreakdash}{\nobreak\hbox{-}}
\setlength{\emergencystretch}{2em}
\multlinegap0pt
\theoremstyle{plain}
\newtheorem{theorem}{Theorem}
\newtheorem{lemma}[theorem]{Lemma}
\theoremstyle{definition}
\newtheorem{definition}[theorem]{Definition}
\DeclareMathOperator{\fes}{fes}
\DeclareMathOperator{\MEG}{meg}
\DeclareMathOperator{\MLN}{mln}
\DeclarePairedDelimiter\abs{\lvert}{\rvert}
\title[The bound in terms of the feedback edge set number]{Monitoring edge-geodetic sets in graphs: the upper bound in terms of the feedback edge set number (Theorem 4.3)}
\author{Subhadeep R. Dev, Sanjana Dey, Florent Foucaud, Krishna Narayanan, Lekshmi Ramasubramony Sulochana}
\date{Source: arXiv:2210.03774v3 (CC BY 4.0)}
\begin{document}
\maketitle
\begin{abstract}
We introduce a new graph-theoretic concept in the area of network monitoring. In this area, one wishes to monitor the vertices and/or the edges of a network (viewed as a graph) in order to detect and prevent failures. Inspired by two notions studied in the literature (edge-geodetic sets and distance-edge-monitoring sets), we define the notion of a monitoring edge-geodetic set (MEG-set for short) of a graph $G$ as an edge-geodetic set $S\subseteq V(G)$ of $G$ (that is, every edge of $G$ lies on some shortest path between two vertices of $S$) with the additional property that for every edge $e$ of $G$, there is a vertex pair $x, y$ of $S$ such that $e$ lies on \emph{all} shortest paths between $x$ and $y$. The motivation is that, if some edge $e$ is removed from the network (for example if it ceases to function), the monitoring probes $x$ and $y$ will detect the failure since the distance between them will increase.

We explore the notion of MEG-sets by deriving the minimum size of a MEG-set for some basic graph classes (trees, cycles, unicyclic graphs, complete graphs, grids, hypercubes, corona products...) and we prove an upper bound using the feedback edge set of the graph.

We also show that determining the smallest size of an MEG-set of a graph is NP-hard, even for graphs of maximum degree at most~9.
\end{abstract}


\noindent\emph{Source and licence.} Monitoring edge-geodetic sets in graphs. Version arXiv:2210.03774v3 (CC BY 4.0), distributed under the Creative Commons Attribution 4.0 International licence, http://creativecommons.org/licenses/by/4.0/, as recorded in the arXiv licence metadata for this exact version and shown on the abstract page https://arxiv.org/abs/2210.03774v3. Full licence notice: licenses/arxiv-2210.03774-CC-BY-4.0.txt.\par\smallskip

\noindent\textit{Editorial scope.} Counted unit: the article's upper bound on the size of a smallest monitoring edge-geodetic set in terms of the feedback edge set number, printed as Theorem 4.3 (source0.tex lines 384--386, source label \texttt{thm:fes}), delivered together with the article's complete original proof of it (source0.tex lines 387--420, one \texttt{proof} environment of 34 lines, adjacent to the statement, with no gap and no deferral). No mathematics is written, repaired or re-derived: statement and proof are the article's own text line for line, in the article's own notation ($G$, $G_b$, $\fes(G)$, $\MEG(G)$, $L(G)$, core vertices, core paths, core cycles, proper core paths, maximal hanging trees, medians $x_1,x_2$, $n_c$, $n_p$), and no line of the retained spans is omitted, substituted or re-flowed. Required context retained uncounted: the article's abstract (lines 67--73, reproduced verbatim as the wrapper's abstract); the article's definition of monitoring edge-geodetic sets and of $\MEG(G)$ (lines 84--91, the definition printed as Definition 1.1); the three preliminary lemmas that the counted proof invokes, namely the simplicial-vertex lemma (lines 139--143, printed as Lemma 2.1) and the cut-vertex lemma (lines 159--164, printed as Lemma 2.3), each with the article's own short proof, and the tree theorem (lines 179--183, printed as Theorem 3.1) with the article's own proof; the statement of the cycle theorem (lines 201--203, printed as Theorem 3.4), whose proof is not retained; the heading of the article's section 4 with its definition of a feedback edge set and of $\fes(G)$ (lines 345--346); the article's terminology paragraph for core vertices, core paths, core cycles, proper core paths, legs, base graph and hanging trees (line 349); the article's sentence introducing its Lemma 4.1 together with that lemma, which the counted proof uses for the count of core vertices and core paths (lines 360--363, printed as Lemma 4.1, itself quoted by the article from the cited references [7] and [15]); the article's Lemma 4.2 with the article's own proof (lines 366--371, printed as Lemma 4.2); and the article's sentence relating the cases $\fes(G)=0$ and $\fes(G)=1$ to the counted bound (line 373). Editorial formatting only, in addition: an AMS \texttt{amsart} preamble that restates the macros and theorem environments the retained lines use -- the operators \texttt{\textbackslash fes}, \texttt{\textbackslash MEG} and \texttt{\textbackslash MLN} exactly as the article's own preamble declares them, the absolute-value command \texttt{\textbackslash abs} that the article obtains from the \texttt{commath} package, and the environments \texttt{theorem}, \texttt{lemma} and \texttt{definition} sharing one counter, as in the article; and sequential numbering of the retained environments, so that they print as Definition 1, Lemma 2, Lemma 3, Theorem 4, Theorem 5, Lemma 6, Lemma 7 and Theorem 8 in order of appearance, whereas the article prints them as Definition 1.1, Lemma 2.1, Lemma 2.3, Theorem 3.1, Theorem 3.4, Lemma 4.1, Lemma 4.2 and Theorem 4.3; the cross-references inside the retained text resolve to the wrapper's numbering. The article is typeset in its own \texttt{article}-class format with its own fonts and colours and the wrapper is an amsart document; the retained \texttt{\textbackslash paragraph} headings print as run-in headings in both. Bibliography: the retained text cites exactly two entries of the article's own in-source \texttt{thebibliography} (lines 623 onward), transcribed verbatim with the article's own printed numbering as their labels ([7] r11 and [15] r12); no other entry of the article's bibliography is reproduced, and no citation outside those two occurs in any retained line. No figure is retained: the article's figure environments (its parameter diagram and its graph illustrations) lie outside every retained span, the retained lines contain no \texttt{\textbackslash includegraphics} command, and the wrapper declares no asset.
\paragraph{Monitoring edge-geodetic sets.}
We now formally define our main concept. 

\begin{definition}
Two vertices $x,y$ \emph{monitor} an edge $e$ in graph $G$ if $e$ belongs to all shortest paths between $x$ and $y$. A set $S$ of vertices of $G$ is called a \emph{monitoring edge-geodetic set} of $G$ (\emph{MEG-set} for short) if, for every edge $e$ of $G$, there is a pair $x,y$ of vertices of $S$ that monitors $e$.
\end{definition}

We denote by $\MEG(G)$ the size of a smallest MEG-set of $G$. We note that $V(G)$ is always an MEG-set of $G$, thus $\MEG(G)$ is always well-defined.

\begin{lemma}\label{simplicial}\label{leaf-node}
In a graph $G$ with at least one edge, any simplicial vertex (in particular, any leaf) belongs to any edge-geodetic set and thus, to any MEG-set of $G$.
\end{lemma}
\begin{proof} Let us consider by contradiction an MEG-set of $G$ that does not contain said simplicial vertex $v$. Any shortest path passing through its neighbors will not pass through $v$, because all the neighbors are adjacent, hence leaving the edges incident to $v$ uncovered, a contradiction.
\end{proof}

\begin{lemma}\label{cutv}
Let $G$ be a graph with a cut-vertex $v$ and $C_1,C_2,\ldots,C_k$ be the $k$ components obtained when removing $v$ from $G$. If $S_1, S_2,\ldots, S_k$ are MEG-sets of the induced subgraphs $G[C_1 \cup \{v\}], G[C_2 \cup \{v\}],\ldots,G[C_k \cup \{v\}]$, then $S = (S_1\cup S_2,\ldots,\cup S_k)\setminus\{v\}$ is an MEG-set of $G$.
\end{lemma}
\begin{proof}
Consider any edge $e$ of $G$, say in $C_1$. Then, there are two vertices $x,y$ of $S_1$ such that $e$ belongs to all shortest paths between $x$ and $y$ in $G_1=G[C_1\cup\{v\}]$. Assume first that $v\notin\{x,y\}$. All shortest paths between $x$ and $y$ in $G$ also exist in $G_1$. Thus, $e$ is monitored by $\{x,y\}\subseteq S$ in $G$. Assume next that $v \in\{x,y\}$: without loss of generality, $v = x$. At least one edge exists in $G[C_2 \cup \{v\}]$, which implies that $S_2\setminus\{v\}$ is nonempty, say, it contains $z$. Then, $e$ is monitored by $y$ and $z$, since $z\in S$. Thus, $S$ monitors all edges of $G$, as claimed.
\end{proof}

\begin{theorem}\label{tt}\label{thm:trees}
For any tree $T$ with at least one edge, the only optimal MEG-set of $T$ consists of the set of leaves of $T$.
\end{theorem}
\begin{proof} The fact that all leaves must be part of any MEG-set follows from Lemma~\ref{leaf-node}, as they are simplicial. For the other side, let $L$ be the set of leaves of $T$. Let $e = xy$ be an edge of $T$ and consider two leaves of $T$, $l_x$ and $l_y$, such that $l_x$ is closer to $x$ than to $y$ and that $l_y$ is closer to $y$ than to $x$. We note that $e$ belongs to the unique (shortest) path between $l_x$ and $l_y$, thus $e$ is monitored by $L$. Hence, $L$ is an MEG-set of $T$. 
\end{proof}

\begin{theorem}\label{cyclet}
Given an $n$-cycle graph $C_n$, for $n = 3$ and $n \geq 5$, $\MEG(C_n) = 3$. Moreover, $\MEG(C_4) = 4$.
\end{theorem}

\section{Relation to feedback edge set number}\label{sec:fes}
A feedback edge set of a graph $G$ is a set of edges which when removed from $G$ leaves a forest. The smallest size of such a feedback edge set of $G$ is denoted by $\fes(G)$ and is sometimes called the \emph{cyclomatic number} of $G$.

We next introduce the following terminology from \cite{r11}. A vertex is a \emph{core vertex} if it has degree at least~3. A path with all internal vertices of degree~2 and whose end-vertices are core vertices is called a \emph{core path}. Do note that we allow the two end-vertices to be equal, but that every other vertex must be distinct. A core path that is a cycle (that is, both end-vertices are equal) is a \emph{core cycle}. For the sake of distinction, a core path that is not a core cycle is called a \emph{proper core path}. We say that a (non-empty) path from a core vertex $u$ to a leaf $v$ is a \textit{leg} of $u$ if all internal vertices of the path have degree~2 ($u$ is not considered to be a part of the leg). The \emph{base graph} of a graph $G$ is the graph of minimum degree~2 obtained from $G$ by iteratively removing vertices of degree~1. A \emph{hanging tree} is a connected subtree of $G$ which is the union of some legs removed from $G$ during the process of creating the base graph $G_b$ of $G$, together with the single vertex of the base graph to which the last removed vertices were adjacent. Thus, $G$ can be decomposed into its base graph and a set of maximal hanging trees. The root of such a maximal hanging tree $T$ is the vertex common to $T$ and $G$.

Based on the aforementioned, we have the following lemma.
\begin{lemma}[\cite{r11,r12}]\label{cores}
Let $G$ be a graph with $\fes(G) = k \geq 2$. The base graph of $G$ has at most $2k$ - $2$ core vertices, that are joined by at most $3k$ - $3$ edge-disjoint core paths. Equivalently, $G$ can be obtained from a multigraph $H$ of order at most $2k-2$ and size at most $3k-3$ by subdividing its edges an arbitrary number of times and iteratively adding degree~1 vertices.
\end{lemma}

\begin{lemma}\label{base}
Let $S$ be an MEG-set of the base graph $G_b$ of $G$ and $L(G)$ be the set of leaves in $G$. Then, $S \cup L(G)$ is an MEG-set of $G$.
\end{lemma}
\begin{proof}
Let $G_b$ be a base graph of $G$. Consider all vertices that are roots of maximal hanging trees on $G_b$. By Theorem~\ref{thm:trees}, the optimal MEG-set of each tree consists of all leaves. We repeatedly apply Lemma~\ref{cutv} to $G$, where for each application of Lemma~\ref{cutv}, the cut-vertex is the root of a hanging tree in consideration.
\end{proof}

Lemma~\ref{leaf-node}, Theorem~\ref{tt} and Lemma~\ref{base} together imply that if $\fes(G) = 0$, then $\MEG(G) \leq \fes(G)+|L(G)|$. Moreover, if $\fes(G)=1$, then $\MEG(G) \leq \fes(G) + |L(G)| + 3$, where $|L(G)|$ is the number of leaves of $G$. We next give a similar bound when $\fes(G)\geq 2$.

\begin{theorem}\label{thm:fes}
If $\fes(G) \geq 2$, then $\MEG(G) \leq 9\fes(G) + |L(G)| - 8$ where $|L(G)|$ is the number of leaves of $G$.
\end{theorem}

\begin{proof}
Let $k = \fes(G)$. We show how to construct a MEG-set $M$ of $G_b$ of order at most $9k-8$ and, by applying Lemma~\ref{base} to $G$, of order $9k-8 + |L(G)|$ for $G$. If an edge $e$ is part of a maximal hanging tree, then by Lemma~\ref{leaf-node} and Lemma~\ref{cutv}, it is monitored by the leaves of $G$ on the maximal hanging tree. $M$ is constructed as follows. 
\begin{itemize}
    \item We let all core vertices of $G_b$ be part of $M$.
    \item One or two internal vertices from each proper core path belongs to $M$, only if the length is at least~2, as explained below.
    \item Two or three internal vertices from each core cycle, as explained below.
\end{itemize}

Consider a proper core path $P$ of length at least~2, with core vertex endpoints $c$ and $c'$, and the median vertex $x_1$ in the case of an even-length path and $x_1$, $x_2$ in the case of an odd-length path, with $d$ edges (on $P$) between the endpoints and the respective medians in $P$. Then, we choose the single median vertex $x_1$ or the two median vertices $x_1$, $x_2$ from each of the core paths and add them to $M$.

For each core cycle $C$, in addition to the core vertex of that cycle, if $C$ has length~4, we add all three non-core vertices of $C$ to $M$. Otherwise, we add two non-core vertices of $C$ to $M$, so that now $C$ has three vertices in $M$. We do this so that these three vertices are as equidistant as possible on the cycle, to be part of $M$ (similar as in Theorem~\ref{cyclet}). % --- note that we may need three vertices if the cycle has length~4, so we conservatively add three vertices in all cases).
%For each core cycle, we conservatively assume the worst case that all core cycles are of length four and include three vertices apart from the core vertex to be part of $M$ (as in Theorem~\ref{cyclet}).

This finishes the description of the construction of $M$.

\medskip

We now show that $M$ monitors all edges of $G_b$. Let $e$ be any edge of $G$. If $e$ lies on a core cycle $C$, assume an origin core vertex of $v_0$. Then, based on Lemma~\ref{cutv} and Theorem~\ref{cyclet}, we deduce that in the worst case, $v_0$ and the two or three other vertices of $C$ in $M$ together suffice to monitor the edges.

If the edge $e$ lies on a proper core path $P$, then let $c$ and $c'$ be the core vertex endpoints of $P$. Let the median vertex of $P$ be $x_1$ in the case of an even-length path and $x_1$, $x_2$ the two medians in the case of an odd-length path. Assume that there are $d$ edges between the medians and closest endpoints of $P$. Without loss of generality, let us say that $e$ lies on the path $P$ such that its closest core vertex is $c$, and closest median $x_1$. 
Given that the distance between $c$ and $x_1$ is $d$ in $P$, the length of any other path between them must be at least $d+1$ (or $d+2$ if $P$ has odd length). Therefore, $c$ and $x_1$ monitor $e$, which lies on the unique shortest path of length~$d$ between them. %(We can similarly argue that if the closest core vertex to $e$ was $c'$ and the closest median vertex was $x_2$, then $c'$ and $x_2$ monitor $e$.) 
If $e$ lies in between the median vertices $x_1$ and $x_2$, then we know that those vertices would monitor $e$ because they are adjacent. %If the path was of even length, then depending on which of the core vertices $c$ and $c'$ was closest to $e$, the distance between the median and the core vertices would be $d$ in $P$ and the length of any other path between them at least $d+1$, ensuring that the median vertex $x_1$ would monitor the edge apart from the core vertices. %Therefore, based on this argument, in the worst case, apart from the core vertices, two median vertices $x_1$ and $x_2$ help monitor the edges on every proper core path $P_i$.
This justifies our construction of $M$, which is an MEG-set of $G_b$, as claimed.

\medskip

Let us now estimate the size of $M$. By Lemma~\ref{cores}, the number of core vertices of $G_b$ is at most $2k-2$, and there are at most $3k-3$ core paths.

If we have core cycles in our graph, then we must note that there can be at most $k$ such cycles in the graph. Indeed, if there were $k+1$ core cycles in the graph, since they are all edge-disjoint, we need at least $k+1$ edges to be removed from $G$ to obtain a forest, a contradiction to the fact that $\fes(G)=k$.

Let $n_c$ be the number of core cycles and $n_p$ be the number of proper core paths. Recall that after adding the core vertices to $M$, we added at most three vertices of each core cycle and two for each core path to $M$. Hence, we have $\abs{M} \leq 3n_c + 2n_p + 2k-2$. Since $n_c \leq k$ and $n_c + n_p \leq 3k-3$ by Lemma~\ref{cores}, we get $\abs{M} \leq 3k + 2(2k-3) + 2k-2 = 9k-8$.

Also note that this value is only reached if all core cycles are of length~4 and all core paths are of odd length.
\end{proof}

\begin{thebibliography}{99}
\bibitem[7]{r11} L. Epstein, A. Levin and G. J. Woeginger. The (weighted) metric dimension of graphs: hard and easy cases. \emph{Algorithmica} 72(4), 1130-1171, 2015.
\bibitem[15]{r12} L. Kellerhals and T. Koana. Parameterized complexity of geodetic set. Proceedings of the 15th International Symposium on Parameterized and Exact Computation (IPEC 2020), \emph{Leibniz International Proceedings in Informatics (LIPIcs)} 180, 20:1--20:14, 2020.
\end{thebibliography}
\end{document}
